% abstract.tex

%TC:break Abstract
\begin{abstract}
Thermal stability limits how far a nanobody can be engineered, yet experimental
melting-temperature (Tm) measurements are scarce. Simulations and physics-based
calculations can label far more sequence variants, but they report quantities
related to Tm rather than Tm itself. We tested these labels as auxiliary supervision
for low-data Tm prediction, with model selection tied only to experimental Tm. On a public benchmark of 57 training, 114 validation, and 396
held-out test sequences, we trained multi-task models in which a Tm predictor
shared a protein-language-model representation with one computational source
task at a time. The source tasks ranged from mutation free energies and
structure-based mutation scores to Rosetta scores for language-model-proposed
variants and molecular-dynamics (MD) trajectory summaries. Alchemical free-energy perturbation (FEP) helped most. It lowered
held-out Tm mean absolute error from 6.61 to 6.26$^\circ$C, a consistent paired
gain of 0.35$^\circ$C (90\% confidence interval, 0.24--0.46$^\circ$C). Enlarging
the ESM2 encoder did not reproduce this gain, and an MD-derived Q-value did not
deliver it at all. Rosetta-scored proposal sets also worked as source labels,
but their edge over random variants stayed within noise. The value of a
computational label therefore depends on whether its physical meaning carries
over to the experimental endpoint, not on how many labels it supplies.
\end{abstract}
%TC:break main

\begin{keywords}
nanobody thermostability | melting temperature | simulation-to-real transfer learning | protein language model | free energy perturbation
\end{keywords}

\begin{corrauthor}
Yasuhiro Matsunaga, ymatsunaga\at riken.jp
\end{corrauthor}
